\chapter{Meshes, models and materials}

Chapter 11 drew everything with \lstinline|DrawCube| and friends. Those are
throwaways: they build geometry from scratch every frame and cannot carry a
texture. This chapter is the step up.

\section{The three layers}

\begin{center}
\begin{tabular}{@{}lp{9.2cm}@{}}
\toprule
\textbf{\lstinline|Mesh|} & raw geometry: vertex positions, normals, texture
coordinates. Lives on the GPU. No colour, no texture. \\[3pt]
\textbf{\lstinline|Material|} & how a surface looks: a shader plus a set of
texture maps (diffuse, normal, roughness\ldots) and a tint. \\[3pt]
\textbf{\lstinline|Model|} & one or more meshes, one or more materials, and a
transform matrix tying them together. \\
\bottomrule
\end{tabular}
\end{center}

A model loaded from a \file{.glb} file arrives with all three filled in. When you
build one yourself, you assemble them:

\begin{lstlisting}
Mesh mesh   = GenMeshCube(2.0f, 2.0f, 2.0f);
Model model = LoadModelFromMesh(mesh);          // wraps it, default white material
\end{lstlisting}

\section{Generating geometry}

You do not need a modelling program to get real meshes:

\begin{center}
\begin{tabular}{@{}ll@{}}
\toprule
\lstinline|GenMeshCube(w, h, l)| & a box \\
\lstinline|GenMeshSphere(radius, rings, slices)| & more rings = smoother, slower \\
\lstinline|GenMeshPlane(w, l, resX, resZ)| & a floor; resolution matters if you deform it \\
\lstinline|GenMeshCylinder(radius, height, slices)| & pillars, trees, barrels \\
\lstinline|GenMeshTorus(radius, size, rad, sides)| & a doughnut \\
\lstinline|GenMeshKnot(...)| & a torus knot, for showing off \\
\lstinline|GenMeshHeightmap(image, size)| & \textbf{terrain from a greyscale image} \\
\lstinline|GenMeshCubicmap(image, size)| & \textbf{a 3D maze from a pixel map} \\
\bottomrule
\end{tabular}
\end{center}

Those last two deserve attention. \lstinline|GenMeshHeightmap| turns a greyscale
picture into hills, where brightness is altitude. \lstinline|GenMeshCubicmap|
turns a black-and-white pixel map into walls --- it is the mini Doom's text map,
except it produces a single mesh instead of hundreds of separate cubes, so it
draws in one call instead of two hundred.

\section{Putting a texture on something}

This is the thing \lstinline|DrawCube| cannot do:

\begin{lstlisting}
Texture2D tex = LoadTexture("wall.png");
model.materials[0].maps[MATERIAL_MAP_DIFFUSE].texture = tex;
\end{lstlisting}

Read it inside out: the model's first material, its diffuse map, its texture.
\lstinline|MATERIAL_MAP_DIFFUSE| is the plain colour map; there are slots for
normal maps, roughness, occlusion and the rest, which only matter once you have a
shader that reads them.

You can tint per draw as well, and the two multiply:

\begin{lstlisting}
DrawModel(model, position, 1.0f, WHITE);    // texture as-is
DrawModel(model, position, 1.0f, RED);      // texture tinted red
\end{lstlisting}

\section{Drawing}

\begin{lstlisting}
DrawModel(model, position, scale, tint);
DrawModelEx(model, position, rotationAxis, rotationAngle, scaleVec3, tint);
DrawModelWires(model, position, scale, tint);
DrawBoundingBox(GetMeshBoundingBox(mesh), GREEN);
\end{lstlisting}

\lstinline|DrawModelEx| is the one you want for anything that turns: an axis, an
angle in degrees, and a per-axis scale.

\begin{lstlisting}
DrawModelEx(model, pos, (Vector3){ 0, 1, 0 }, angle, (Vector3){ 1, 1, 1 }, WHITE);
\end{lstlisting}

Rotating around \lstinline|(0,1,0)| means spinning about the vertical axis, which
is what you want for a character turning or a pickup idling.

\section{Loading real models}

\begin{lstlisting}
Model robot = LoadModel("robot.glb");       // once
DrawModel(robot, position, 1.0f, WHITE);    // every frame
UnloadModel(robot);                         // once
\end{lstlisting}

raylib reads \file{.obj}, \file{.gltf}, \file{.glb}, \file{.iqm}, \file{.m3d} and
\file{.vox}. Prefer \textbf{\file{.glb}}: it packs the mesh, the materials and the
textures into one file, so there is nothing to go missing and no paths to fix.

\gotcha{\lstinline|UnloadModel| frees the meshes it owns but deliberately does
\emph{not} free textures or shaders --- you might be sharing them between models,
so raylib leaves them to you. Unload your textures yourself, exactly once.}

\section{Animation}

\begin{lstlisting}
int animCount = 0;
ModelAnimation *anims = LoadModelAnimations("robot.glb", &animCount);
int frame = 0;

// in the loop:
frame = (frame + 1) % anims[0].frameCount;
UpdateModelAnimation(robot, anims[0], frame);
DrawModel(robot, position, 1.0f, WHITE);

UnloadModelAnimations(anims, animCount);
\end{lstlisting}

The model must have been exported with a skeleton, \file{.glb} and \file{.iqm}
being the usual choices. Note that stepping \lstinline|frame| once per frame ties
the animation to your frame rate --- use a timer, as in chapter 5, if you care.

\section{Billboards: the old trick, still useful}

A billboard is a flat image that always turns to face the camera. Every enemy in
Doom was one.

\begin{lstlisting}
DrawBillboard(camera, texture, position, size, WHITE);
DrawBillboardRec(camera, texture, sourceRec, position, size, WHITE);  // one sprite of a sheet
\end{lstlisting}

They are cheap, they look deliberately retro, and \lstinline|DrawBillboardRec|
lets you animate them from a sprite sheet exactly as in chapter 7. For particles,
smoke, flames and pickups they are still the right answer today.

\section{When to stop using DrawCube}

\begin{center}
\begin{tabular}{@{}ll@{}}
\toprule
\textbf{Stay with immediate shapes if} & \textbf{Move to models when} \\
\midrule
prototyping & you want textures \\
a handful of objects & you want lighting via a shader \\
debug visualisation & you draw hundreds of the same thing \\
flat colour is fine & you need animation \\
\bottomrule
\end{tabular}
\end{center}

And when you are drawing the same mesh thousands of times --- grass, crowds,
asteroids --- there is \lstinline|DrawMeshInstanced|, which sends one mesh and an
array of transforms in a single call. It is the difference between 50 FPS and
2000.

\tryit{Module 8 generates a cube, a sphere and a torus, puts a generated
checkerboard on the cube's material, draws a billboard, and lets you click any
object to pick it with a ray. TAB switches between models, immediate shapes and
wireframe so you can compare them.}{./cheatsheet/build.sh 8}
